\documentclass[11pt]{article}
\usepackage[margin=1in]{geometry}
\title{VTA, Lesioned Hemisphere --- Interpretation}
\author{GSE233866, 6-OHDA lesion cohort}
\date{}

\begin{document}
\maketitle

VTA loses proportionally less of its population to lesioning than SNc does
(4{,}426 $\to$ 1{,}720 cells, 2.6$\times$, versus SNc's 3.8$\times$) ---
consistent with this dataset's premise that SNc is the more vulnerable
region. The resulting 1{,}720-cell network selects soft power 6
($R^2=0.939$), matching the intact-SNc network's power despite the different
cell population.

CACNA1D lands in the \textbf{blue} module with $\mathrm{kME}=0.283$ ---
weaker than either SNc lesion-arm value (intact: 0.346, lesioned: 0.398),
placing lesioned-VTA closer to the healthy-baseline connectivity range for
this gene. This network is the standalone counterpart to the formal
region-comparison test in \texttt{snc\_vs\_vta\_combined\_dme\_lesioned/},
which is where the statistically-supported SNc-vs-VTA claim during
lesioning is actually established.

\end{document}
